\documentclass{article}
\usepackage[utf8]{inputenc}
\usepackage{amsmath}
\usepackage{geometry}
\usepackage{amsfonts}
\geometry{a4paper, margin=1in}

\title{$\boxed{\textbf{FORMULARIO (VERSIÓN EXTENDIDA)}}$}

\begin{document}
	\maketitle

	\section*{Constantes}
	\addcontentsline{toc}{subsection}{Constantes}
		\begin{itemize}
			\item \textbf{(0.1)} $\boxed{\epsilon_{0}=8,854\cdot10^{-12}\,\text{F/m}}$ Permitividad eléctrica del vacío. Resistencia que ofrece el vacío a la formación de un campo eléctrico.
			\item \textbf{(0.2)} $\boxed{\frac{1}{4\pi\epsilon_{0}} \approx 9\cdot10^{9}\,\text{N}\cdot\text{m}^{2}/\text{C}^{2}}$ Constante de Coulomb ($K$).
			\item \textbf{(0.3)} $\boxed{\mu_{0}=4\pi\times10^{-7}\,\text{H/m}}$ Permeabilidad magnética del vacío. Capacidad del vacío para dejar pasar líneas de campo magnético.
			\item \textbf{(0.4)} $\boxed{c \approx 3\cdot10^{8}\,\text{m/s}}$ Velocidad de la luz en el vacío.
		\end{itemize}

	\section*{Áreas y volúmenes}
	\addcontentsline{toc}{subsection}{Áreas y volúmenes}
		\begin{itemize}
			\item \textbf{(0.5)} $\boxed{L=2\pi r}$ Longitud de la circunferencia.
				\begin{itemize}
					\item $r$: Radio de la circunferencia.
				\end{itemize}
			\item \textbf{(0.6)} $\boxed{S=\pi r^{2}}$ Área del círculo.
			\item \textbf{(0.7)} $\boxed{S=4\pi r^{2}}$ Área superficial de una esfera.
			\item \textbf{(0.8)} $\boxed{V=\frac{4}{3}\pi r^{3}}$ Volumen de una esfera.
			\item \textbf{(0.9)} $\boxed{S=2\pi r^{2}+2\pi rL}$ Área superficial de un cilindro.
				\begin{itemize}
					\item $r$: Radio de la base del cilindro.
					\item $L$: Altura (o longitud) del cilindro.
				\end{itemize}
			\item \textbf{(0.10)} $\boxed{V=\pi r^{2}L}$ Volumen de un cilindro.
		\end{itemize}

	\section{Óptica y fotónica}
		\subsection*{Ondas}
			\begin{itemize}
				\item \textbf{(1.1)} $\boxed{f(x,t)=A\,\sin(kx-\omega t+\phi)}$ Función de la propagación de una onda armónica en el espacio $x$ y en el tiempo $t$.
					\begin{itemize}
						\item $f(x,t)$: Función de onda. Representa el desplazamiento de la onda en una posición $x$ y un instante de tiempo $t$. Se mide en metros ($m$).
						\item $A$: Amplitud. Es el valor máximo de la oscilación. Indica la intensidad; es decir, lo lejos que llega la onda desde su punto de equilibrio. Se mide en metros ($m$).
						\item $k$: Número de onda. Indica cómo de rápido cambia la onda en el espacio. Se relaciona con la longitud de onda ($\lambda$) mediante la fórmula $k = \frac{2\pi}{\lambda}$. Se mide en radianes por metro ($rad/m$).
						\item $x$: Posición espacial. Es la coordenada donde se evalúa la onda. Se mide en metros ($m$).
						\item $\omega$: Frecuencia angular. Indica cómo de rápido cambia la onda en el tiempo. Se relaciona con el periodo ($T$) y la frecuencia ($f$) mediante $\omega = \frac{2\pi}{T} = 2\pi f$. Se mide en radianes por segundo ($rad/s$).
						\item $t$: Tiempo. Es el instante en el que se evalúa la onda. Se mide en segundos ($s$).
						\item $\phi$: Constante de fase. Determina el estado inicial de la oscilación en el tiempo $t=0$ y la posición $x=0$. Indica el "desplazamiento" inicial de la onda. Se mide en radianes ($rad$).
					\end{itemize}

				\item \textbf{(1.2)} $\boxed{\lambda=\frac{2\pi}{k}}$ Longitud de onda. Distancia entre dos puntos consecutivos que están en la misma fase.
					\begin{itemize}
						\item $\lambda$: Longitud de onda. Distancia entre dos puntos consecutivos que están en la misma fase. Se mide en metros ($m$).
						\item $k$: Número de onda. Indica cómo de rápido cambia la onda en el espacio. Se relaciona con la longitud de onda ($\lambda$) mediante la fórmula $k = \frac{2\pi}{\lambda}$. Se mide en radianes por metro ($rad/m$).
					\end{itemize}

				\item \textbf{(1.3a)} $\boxed{T=\frac{1}{f}}$ Periodo de la onda a partir de la frecuencia. Tiempo que tarda en completarse un ciclo.
					\begin{itemize}
						\item $T$: Periodo de la onda. Tiempo que tarda en completarse un ciclo. Se mide en segundos ($s$).
						\item $f$: Frecuencia de onda. Número de ciclos por segundo. Se mide en Hercios ($Hz$).
					\end{itemize}

				\item \textbf{(1.3b)} $\boxed{T=\frac{2\pi}{\omega}}$ Periodo de la onda a partir de la frecuencia angular. Tiempo que tarda en completarse un ciclo.
					\begin{itemize}
						\item $T$: Periodo de la onda. Tiempo que tarda en completarse un ciclo. Se mide en segundos ($s$).
						\item $\omega$: Frecuencia angular. Indica cómo de rápido cambia la onda en el tiempo. Se relaciona con el periodo ($T$) y la frecuencia ($f$) mediante $\omega = \frac{2\pi}{T} = 2\pi f$. Se mide en radianes por segundo ($rad/s$).
					\end{itemize}

				\item \textbf{(1.4a)} $\boxed{v=\lambda f}$ Velocidad de propagación de la onda a partir de la longitud de onda y la frecuencia.
					\begin{itemize}
						\item $v$: Velocidad de propagación de la onda. Rapidez con la que se desplaza la onda a través del medio. Se mide en metros por segundo ($m/s$).
						\item $\lambda$: Longitud de onda. Distancia entre dos puntos consecutivos que están en la misma fase. Se mide en metros ($m$).
						\item $f$: Frecuencia de onda. Número de ciclos por segundo. Se mide en Hercios ($Hz$).
					\end{itemize}

				\item \textbf{(1.4b)} $\boxed{v=\frac{\omega}{k}}$ Velocidad de propagación de la onda a partir de la frecuencia angular y el número de onda.
					\begin{itemize}
						\item $v$: Velocidad de propagación de la onda. Rapidez con la que se desplaza la onda a través del medio. Se mide en metros por segundo ($m/s$).
						\item $\omega$: Frecuencia angular. Indica cómo de rápido cambia la onda en el tiempo. Se relaciona con el periodo ($T$) y la frecuencia ($f$) mediante $\omega = \frac{2\pi}{T} = 2\pi f$. Se mide en radianes por segundo ($rad/s$).
						\item $k$: Número de onda. Indica cómo de rápido cambia la onda en el espacio. Se relaciona con la longitud de onda ($\lambda$) mediante la fórmula $k = \frac{2\pi}{\lambda}$. Se mide en radianes por metro ($rad/m$).
					\end{itemize}
			\end{itemize}

		\subsection*{Óptica geométrica}
			\begin{itemize}
				\item \textbf{(1.5)} $\boxed{n=\frac{v_{0}}{v}}$ Índice de refracción. Es una medida que indica cuánto se reduce la velocidad de la luz al pasar del vacío a un medio material.
					\begin{itemize}
						\item $n$: Índice de refracción. Es una medida de cuánto reduce la luz su velocidad en un medio. Es adimensional (no tiene unidades).
						\item $v_0$: Velocidad de la luz en el vacío ($\approx 3 \times 10^8 \, m/s$). Se mide en metros por segundo ($m/s$).
						\item $v$: Velocidad de propagación de la luz en el medio material. Se mide en metros por segundo ($m/s$).
					\end{itemize}

				\item \textbf{(1.6)} $\boxed{n_{1}\sin\,\theta_{1}=n_{2}\sin\,\theta_{2}}$ Ley de Snell para la refracción. Relaciona los ángulos de incidencia y refracción con los índices de refracción de los medios al atravesar una superficie de separación.
					\begin{itemize}
						\item $n_1$: Índice de refracción del medio 1 (medio de incidencia). Es una medida de cuánto reduce la luz su velocidad en el primer medio. Es adimensional.
						\item $\theta_1$: Ángulo de incidencia. Es el ángulo formado entre el rayo de luz incidente y la normal (línea perpendicular a la superficie de separación). Se mide en grados ($^\circ$) o radianes ($rad$).
						\item $n_2$: Índice de refracción del medio 2 (medio de refracción). Es una medida de cuánto reduce la luz su velocidad en el segundo medio. Es adimensional.
						\item $\theta_2$: Ángulo de refracción. Ángulo formado entre el rayo refractado y la normal. Se mide en grados ($^\circ$) o radianes ($rad$).
					\end{itemize}

					Conceptos clave para interpretar esta fórmula:
					\begin{itemize}
						\item La Normal: Los ángulos $\theta_1$ y $\theta_2$ se miden respecto a la \textbf{línea perpendicular a la superficie de contacto}, no respecto a la superficie de contacto.
						\item Si $n_2 > n_1$, la luz se acerca a la normal; el rayo "se dobla" hacia adentro.
						\item Si $n_2 < n_1$, la luz se aleja de la normal; el rayo "se dobla" hacia afuera.
					\end{itemize}

				\item \textbf{(1.7)} $\boxed{\sin\,\theta_{1}=\sqrt{(\frac{n_{2}}{n_{1}})^{2}-(\frac{n_{3}}{n_{1}})^{2}}}$ Ángulo límite / apertura numérica. Determina el cono de aceptación o el ángulo máximo de incidencia para que la luz quede confinada mediante reflexión interna total (por ejemplo, dentro de una fibra óptica).
					\begin{itemize}
						\item $\theta_1$: Ángulo límite o de apertura (ángulo de aceptación). Es el ángulo máximo de incidencia para el cual la luz se transmitirá de forma confinada. Se mide en grados ($^\circ$) o radianes ($rad$).
						\item $n_1$: Índice de refracción del medio desde el cual incide la luz (ej. núcleo de fibra óptica). Es adimensional.
						\item $n_2$: Índice de refracción del segundo medio (ej. revestimiento). Es adimensional.
						\item $n_3$: Índice de refracción del medio externo (ej. aire). Es adimensional.
					\end{itemize}
			\end{itemize}

		\subsection*{Óptica ondulatoria y electromagnetismo}
			\begin{itemize}
				\item \textbf{(1.8)} $\boxed{\vert{}\vec{E}\vert{}=v_{0}\vert{}\vec{B}\vert{}}$ Relación de amplitudes entre campo eléctrico y magnético.
					\begin{itemize}
						\item $\vert{}\vec{E}\vert{}$: Módulo del campo eléctrico instantáneo. Se mide en Voltios por metro ($V/m$) o Newtons por Culombio ($N/C$).
						\item $v_0$: Velocidad de propagación de la onda en el medio. Si la onda viaja por el vacío, $v_0$ es exactamente igual a $c$ (la velocidad de la luz, $\approx 3 \times 10^8 \, m/s$). Se mide en metros por segundo ($m/s$).
						\item $\vert{}\vec{B}\vert{}$: Módulo del campo magnético instantáneo. Se mide en Teslas ($T$).
					\end{itemize}

				\item \textbf{(1.9)} $\boxed{v_{0}=\frac{1}{\sqrt{\mu_{0}\epsilon_{0}}}}$ Velocidad de la onda electromagnética en el vacío.
					\begin{itemize}
						\item $v_0$: Velocidad de propagación de la onda en el medio. Si la onda viaja por el vacío, $v_0$ es exactamente igual a $c$ (la velocidad de la luz, $\approx 3 \times 10^8 \, m/s$). Se mide en metros por segundo ($m/s$).
						\item $\mu_0$: Permeabilidad magnética del vacío. Constante física que indica cómo el vacío responde a un campo magnético. Su valor es $4\pi \times 10^{-7} \, T\cdot m/A$.
						\item $\epsilon_{0}$: Permitividad eléctrica del vacío. Constante física cuyo valor es $8,854\cdot10^{-12}\,\text{F/m}$. Describe la resistencia que ofrece el vacío a la formación de un campo eléctrico.
					\end{itemize}

				\item \textbf{(1.10)} $\boxed{\vec{E}(\vec{r},t)=\vec{E}_{0}\sin(\vec{k}\vec{r}-\omega t+\phi)}$ Expresión vectorial del campo eléctrico.
					\begin{itemize}
						\item $\vec{E}(\vec{r},t)$: Campo eléctrico vectorial en función de la posición y el tiempo. Se mide en Voltios por metro ($V/m$) o Newtons por Culombio ($N/C$).
						\item $\vec{E}_0$: Vector amplitud. Indica la dirección de oscilación del campo eléctrico y su valor máximo. Se mide en Voltios por metro ($V/m$) o Newtons por Culombio ($N/C$).
						\item $\vec{k}$: Vector número de onda. Indica la dirección de propagación de la onda. Se mide en radianes por metro ($rad/m$).
						\item $\vec{r}$: Vector posición espacial $(x,y,z)$. Indica el punto del espacio donde se evalúa el campo. Se mide en metros ($m$).
						\item $\omega$: Frecuencia angular. Indica cómo de rápido cambia la onda en el tiempo. Se mide en radianes por segundo ($rad/s$).
						\item $t$: Tiempo. Es el instante en el que se evalúa la onda. Se mide en segundos ($s$).
						\item $\phi$: Constante de fase. Determina el estado inicial de la oscilación. Se mide en radianes ($rad$).
					\end{itemize}

				\item \textbf{(1.11)} $\boxed{u_{e}=\frac{1}{2}\epsilon_{0}E^{2}}$ Densidad de energía eléctrica instantánea.
					\begin{itemize}
						\item $u_e$: Densidad de energía eléctrica instantánea. Energía almacenada por el campo eléctrico en un volumen dado. Se mide en Julios por metro cúbico ($J/m^3$).
						\item $\epsilon_{0}$: Permitividad eléctrica del vacío. Constante física cuyo valor es $8,854\cdot10^{-12}\,\text{F/m}$. Describe la resistencia que ofrece el vacío a la formación de un campo eléctrico.
						\item $E$: Módulo del campo eléctrico instantáneo. Se mide en Voltios por metro ($V/m$) o Newtons por Culombio ($N/C$).
					\end{itemize}

				\item \textbf{(1.12)} $\boxed{u_{m}=\frac{B^{2}}{2\mu_{0}}}$ Densidad de energía magnética instantánea.
					\begin{itemize}
						\item $u_m$: Densidad de energía magnética instantánea. Energía almacenada por el campo magnético en un volumen dado. Se mide en Julios por metro cúbico ($J/m^3$).
						\item $B$: Módulo del campo magnético instantáneo. Se mide en Teslas ($T$).
						\item $\mu_0$: Permeabilidad magnética del vacío. Constante física que indica cómo el vacío responde a un campo magnético. Su valor es $4\pi \times 10^{-7} \, T\cdot m/A$.
					\end{itemize}

				\item \textbf{(1.13)} $\boxed{u=\epsilon_{0}E^{2}}$ Densidad de energía total instantánea.
					\begin{itemize}
						\item $u$: Densidad de energía total instantánea. Suma de las energías eléctrica y magnética. Se mide en Julios por metro cúbico ($J/m^3$).
						\item $\epsilon_{0}$: Permitividad eléctrica del vacío. Constante física cuyo valor es $8,854\cdot10^{-12}\,\text{F/m}$. Describe la resistencia que ofrece el vacío a la formación de un campo eléctrico.
						\item $E$: Módulo del campo eléctrico instantáneo. Se mide en Voltios por metro ($V/m$) o Newtons por Culombio ($N/C$).
					\end{itemize}

				\item \textbf{(1.14)} $\boxed{\langle u\rangle=\frac{E_{0}B_{0}}{2\mu_{0}v_{0}}}$ Densidad de energía media. Calcula cuánta energía electromagnética hay almacenada, en promedio, en un metro cúbico de espacio mientras la onda viaja a través de él.
					\begin{itemize}
						\item $\langle u\rangle$: Densidad de energía media. Se mide en Julios por metro cúbico ($J/m^3$). Los corchetes angulares $\langle \rangle$ en física denotan un valor promedio.
						\item $E_0$: Amplitud del campo eléctrico. Es el valor máximo que alcanza el campo eléctrico de la onda. Se mide en Voltios por metro ($V/m$) o Newtons por Culombio ($N/C$).
						\item $B_0$: Amplitud del campo magnético. Es el valor máximo que alcanza el campo magnético de la onda. Se mide en Teslas ($T$).
						\item $\mu_0$: Permeabilidad magnética del vacío. Constante física que indica cómo el vacío responde a un campo magnético. Su valor es $4\pi \times 10^{-7} \, T\cdot m/A$.
						\item $v_0$: Velocidad de propagación de la onda en el medio. Si la onda viaja por el vacío, $v_0$ es exactamente igual a $c$ (la velocidad de la luz, $\approx 3 \times 10^8 \, m/s$). Se mide en metros por segundo ($m/s$).
					\end{itemize}

				\item \textbf{(1.15)} $\boxed{I=\langle u\rangle v_{0}}$ Relación entre intensidad y densidad. La intensidad de una onda es su densidad de energía multiplicada por la velocidad a la que se mueve ese volumen.
					\begin{itemize}
						\item $I$: Intensidad de la onda. Potencia promedio transferida a través de una unidad de área perpendicular a la dirección de propagación. Se mide en Vatios por metro cuadrado ($W/m^2$).
						\item $\langle u\rangle$: Densidad de energía media. Se mide en Julios por metro cúbico ($J/m^3$). Los corchetes angulares $\langle \rangle$ en física denotan un valor promedio.
						\item $v_0$: Velocidad de propagación de la onda en el medio. Si la onda viaja por el vacío, $v_0$ es exactamente igual a $c$ (la velocidad de la luz, $\approx 3 \times 10^8 \, m/s$). Se mide en metros por segundo ($m/s$).
					\end{itemize}

				\item \textbf{(1.16)} $\boxed{I=\frac{E_{0}B_{0}}{2\mu_{0}}}$ Intensidad a partir de los campos. Es la forma más directa de calcular la intensidad si se conocen los valores máximos de los campos eléctrico y magnético.
					\begin{itemize}
						\item $I$: Intensidad de la onda. Potencia promedio transferida a través de una unidad de área perpendicular a la dirección de propagación. Se mide en Vatios por metro cuadrado ($W/m^2$).
						\item $E_0$: Amplitud del campo eléctrico. Es el valor máximo que alcanza el campo eléctrico de la onda. Se mide en Voltios por metro ($V/m$) o Newtons por Culombio ($N/C$).
						\item $B_0$: Amplitud del campo magnético. Es el valor máximo que alcanza el campo magnético de la onda. Se mide en Teslas ($T$).
						\item $\mu_0$: Permeabilidad magnética del vacío. Constante física que indica cómo el vacío responde a un campo magnético. Su valor es $4\pi \times 10^{-7} \, T\cdot m/A$.
					\end{itemize}

				\item \textbf{(1.17)} $\boxed{\vec{S}=\frac{\vec{E}\times\vec{B}}{\mu_{0}}}$ Vector de Poynting.
					\begin{itemize}
						\item $\vec{S}$: Vector de Poynting. Indica la dirección de propagación de la energía y su magnitud es la intensidad instantánea. Es perpendicular tanto a $\vec{E}$ como a $\vec{B}$. Se mide en Vatios por metro cuadrado ($W/m^2$).
						\item $\vec{E}$: Vector campo eléctrico instantáneo. Se mide en Voltios por metro ($V/m$) o Newtons por Culombio ($N/C$).
						\item $\vec{B}$: Vector campo magnético instantáneo. Se mide en Teslas ($T$).
						\item $\mu_0$: Permeabilidad magnética del vacío. Constante física que indica cómo el vacío responde a un campo magnético. Su valor es $4\pi \times 10^{-7} \, T\cdot m/A$.
					\end{itemize}

				\item \textbf{(1.18)} $\boxed{P_{r}=\frac{I}{v_{0}}}$ Presión de radiación. Fuerza por unidad de área ejercida por la onda al chocar contra una superficie.
					\begin{itemize}
						\item $P_r$: Presión de radiación. Fuerza por unidad de área ejercida por la onda al chocar contra una superficie. Se mide en Pascales ($Pa$) o Newtons por metro cuadrado ($N/m^2$).
						\item $I$: Intensidad de la onda. Potencia promedio transferida a través de una unidad de área perpendicular a la dirección de propagación. Se mide en Vatios por metro cuadrado ($W/m^2$).
						\item $v_0$: Velocidad de propagación de la onda en el medio. Si la onda viaja por el vacío, $v_0$ es exactamente igual a $c$ (la velocidad de la luz, $\approx 3 \times 10^8 \, m/s$). Se mide en metros por segundo ($m/s$).
					\end{itemize}
			\end{itemize}

		\subsection*{Difracción e interferencia}
			\begin{itemize}
				\item \textbf{(1.19)} $\boxed{I=I_{1}+I_{2}+2\sqrt{I_{1}I_{2}}\cos\,\delta}$ Intensidad resultante por interferencia de dos ondas.
					\begin{itemize}
						\item $I$: Intensidad de la onda. Potencia promedio transferida a través de una unidad de área perpendicular a la dirección de propagación. Se mide en Vatios por metro cuadrado ($W/m^2$).
						\item $I_1, I_2$: Intensidad de la onda 1 y 2. Potencia promedio transferida a través de una unidad de área perpendicular a la dirección de propagación. Se mide en Vatios por metro cuadrado ($W/m^2$).
						\item $\delta$: Diferencia de fase (desfase) entre las dos ondas cuando se superponen. Se mide en radianes ($rad$).
					\end{itemize}

				\item \textbf{(1.20)} $\boxed{\delta=k(r_{1}-r_{2})=k\Delta r}$ Desfase en función de la diferencia de caminos.
					\begin{itemize}
						\item $\delta$: Diferencia de fase (desfase). Se mide en radianes ($rad$).
						\item $k$: Número de onda. Indica cómo de rápido cambia la onda en el espacio. Se mide en radianes por metro ($rad/m$).
						\item $r_1, r_2$: Distancia geométrica recorrida por la onda 1 y 2 hasta el punto de encuentro. Se mide en metros ($m$).
						\item $\Delta r$: Diferencia geométrica de las distancias recorridas por ambas ondas (diferencia de camino óptico). Se mide en metros ($m$).
					\end{itemize}

				\item \textbf{(1.21)} $\boxed{k=\frac{2\pi}{\lambda}}$ Número de onda.
					\begin{itemize}
						\item $k$: Número de onda. Indica cómo de rápido cambia la onda en el espacio. Se mide en radianes por metro ($rad/m$).
						\item $\lambda$: Longitud de onda. Distancia entre dos puntos consecutivos que están en la misma fase. Se mide en metros ($m$).
					\end{itemize}

				\item \textbf{(1.22)} $\boxed{\delta=\frac{2\pi}{\lambda}\Delta r}$ Desfase explícito según la longitud de onda.
					\begin{itemize}
						\item $\delta$: Diferencia de fase (desfase). Se mide en radianes ($rad$).
						\item $\lambda$: Longitud de onda. Distancia entre dos puntos consecutivos que están en la misma fase. Se mide en metros ($m$).
						\item $\Delta r$: Diferencia geométrica de las distancias recorridas por ambas ondas. Se mide en metros ($m$).
					\end{itemize}
			\end{itemize}

		\subsection*{Doble rendija}
			\begin{itemize}
				\item \textbf{(1.23)} $\boxed{\Delta r=d\sin\,\theta}$ Diferencia de caminos en el experimento de doble rendija de Young.
					\begin{itemize}
						\item $\Delta r$: Diferencia geométrica de las distancias recorridas por ambas ondas. Se mide en metros ($m$).
						\item $d$: Distancia de separación entre las dos rendijas. Se mide en metros ($m$).
						\item $\theta$: Ángulo de desviación respecto a la franja central. Se mide en grados ($^\circ$) o radianes ($rad$).
					\end{itemize}    

				\item \textbf{(1.24)} $\boxed{\delta=\frac{2\pi}{\lambda}d\sin\,\theta}$ Desfase en función de las variables de la doble rendija.
					\begin{itemize}
						\item $\delta$: Diferencia de fase (desfase). Se mide en radianes ($rad$).
						\item $\lambda$: Longitud de onda. Se mide en metros ($m$).
						\item $d$: Distancia de separación entre las dos rendijas. Se mide en metros ($m$).
						\item $\theta$: Ángulo de desviación respecto a la franja central. Se mide en grados ($^\circ$) o radianes ($rad$).
					\end{itemize}

				\item \textbf{(1.25)} $\boxed{d\sin\,\theta=m\lambda, \quad m=0,1,2,3\dots}$ Condición de Máximos (interferencia constructiva).
					\begin{itemize}
						\item $d$: Distancia de separación entre las dos rendijas. Se mide en metros ($m$).
						\item $\theta$: Ángulo de desviación respecto a la franja central. Se mide en grados ($^\circ$) o radianes ($rad$).
						\item $m$: Orden de interferencia. Es un número entero utilizado para ubicar las franjas brillantes. Es adimensional.
						\item $\lambda$: Longitud de onda. Se mide en metros ($m$).
					\end{itemize}

				\item \textbf{(1.26)} $\boxed{d\sin\,\theta=(m-1/2)\lambda, \quad m=1,2,3\dots}$ Condición de Mínimos (interferencia destructiva).
					\begin{itemize}
						\item Iguales a la fórmula (1.25), pero $m$ comienza en 1 y representa las franjas oscuras.
					\end{itemize}

				\item \textbf{(1.27)} $\boxed{\tan\,\theta=\frac{y}{L}}$ Relación trigonométrica para proyectar la desviación en una pantalla.
					\begin{itemize}
						\item $\theta$: Ángulo de desviación. Se mide en grados ($^\circ$) o radianes ($rad$).
						\item $y$: Distancia vertical en la pantalla desde la franja central (origen) hasta la franja observada. Se mide en metros ($m$).
						\item $L$: Distancia en línea recta desde las rendijas hasta la pantalla. Se mide en metros ($m$).
					\end{itemize}

				\item \textbf{(1.28)} $\boxed{y=m\frac{\lambda L}{d}, \quad m=0,1,2,3\dots}$ Ubicación lineal de los Máximos en la pantalla.
					\begin{itemize}
						\item $y$: Distancia de la franja brillante al centro. Se mide en metros ($m$).
						\item $m$: Orden del máximo. Es adimensional.
						\item $\lambda$: Longitud de onda. Se mide en metros ($m$).
						\item $L$: Distancia a la pantalla. Se mide en metros ($m$).
						\item $d$: Distancia entre rendijas. Se mide en metros ($m$).
					\end{itemize}

				\item \textbf{(1.29)} $\boxed{y=(m-1/2)\frac{\lambda L}{d}, \quad m=1,2,3\dots}$ Ubicación lineal de los Mínimos en la pantalla.
					\begin{itemize}
						\item $y$: Distancia de la franja oscura al centro. Se mide en metros ($m$).
						\item Resto de variables: Idénticas a la fórmula (1.28).
					\end{itemize}
			\end{itemize}

		\subsection*{Única abertura}
			\begin{itemize}
				\item \textbf{(1.30)} $\boxed{a\sin\,\theta=m\lambda, \quad m=1,2,3\dots}$ Condición de Mínimos de difracción de una rendija simple.
				\begin{itemize}
					\item $a$: Ancho de la rendija. Se mide en metros ($m$).
					\item $\theta$: Ángulo de desviación hacia la zona de oscuridad. Se mide en grados ($^\circ$) o radianes ($rad$).
					\item $m$: Orden del mínimo de difracción. Es adimensional.
					\item $\lambda$: Longitud de onda. Se mide en metros ($m$).
				\end{itemize}
			\end{itemize}

		\subsection*{Óptica cuántica}
			\begin{itemize}
				\item \textbf{(1.31)} $\boxed{E=hf}$ Energía de un fotón cuantizado.
					\begin{itemize}
						\item $E$: Energía del fotón. Se mide en Julios ($J$) o Electrón-voltios ($eV$).
						\item $h$: Constante de Planck. Se mide en Julios-segundo ($J\cdot s$).
						\item $f$: Frecuencia de la onda. Se mide en Hercios ($Hz$).
					\end{itemize}    

				\item \textbf{(1.32)} $\boxed{h=6,626\cdot10^{-34}\,\text{J}\cdot\text{s}}$ Constante de Planck.
					\begin{itemize}
						\item $h$: Constante fundamental de la física cuántica que relaciona la energía de una partícula con su frecuencia.
					\end{itemize}

				\item \textbf{(1.33)} $\boxed{\lambda=\frac{h}{mv}}$ Longitud de onda de De Broglie para partículas.
					\begin{itemize}
						\item $\lambda$: Longitud de onda asociada a la partícula. Se mide en metros ($m$).
						\item $h$: Constante de Planck. Se mide en Julios-segundo ($J\cdot s$).
						\item $m$: Masa de la partícula. Se mide en kilogramos ($kg$).
						\item $v$: Velocidad de la partícula. Se mide en metros por segundo ($m/s$).
					\end{itemize}

				\item \textbf{(1.34)} $\boxed{f=\frac{E_{2}-E_{1}}{h}}$ Frecuencia del fotón en un salto cuántico.
					\begin{itemize}
						\item $f$: Frecuencia del fotón emitido o absorbido. Se mide en Hercios ($Hz$).
						\item $E_2$: Energía del estado de mayor energía (excitado). Se mide en Julios ($J$) o Electrón-voltios ($eV$).
						\item $E_1$: Energía del estado de menor energía (fundamental). Se mide en Julios ($J$) o Electrón-voltios ($eV$).
						\item $h$: Constante de Planck. Se mide en Julios-segundo ($J\cdot s$).
					\end{itemize}

				\item \textbf{(1.35)} $\boxed{\frac{dn_{1}}{dt}=-B_{12}\rho n_{1}}$ Absorción estimulada (ecuaciones de tasa de Einstein).
					\begin{itemize}
						\item $\frac{dn_1}{dt}$: Tasa de variación de la población en el nivel inferior respecto al tiempo.
						\item $B_{12}$: Coeficiente de probabilidad de absorción estimulada de Einstein.
						\item $\rho$: Densidad de energía de la radiación incidente. Se mide en Julios por metro cúbico por Hercio ($J/(m^3\cdot Hz)$).
						\item $n_1$: Número de átomos en el estado fundamental.
					\end{itemize}

				\item \textbf{(1.36)} $\boxed{\frac{dn_{2}}{dt}=-B_{21}\rho n_{2}}$ Emisión estimulada (ecuaciones de tasa de Einstein).
					\begin{itemize}
						\item $\frac{dn_2}{dt}$: Tasa de variación de la población en el nivel superior debido a estímulos externos.
						\item $B_{21}$: Coeficiente de probabilidad de emisión estimulada de Einstein.
						\item $\rho$: Densidad de energía de la radiación incidente.
						\item $n_2$: Número de átomos en el estado excitado.
					\end{itemize}

				\item \textbf{(1.37)} $\boxed{\frac{dn_{2}}{dt}=-A_{21}n_{2}}$ Emisión espontánea (ecuaciones de tasa de Einstein).
					\begin{itemize}
						\item $\frac{dn_2}{dt}$: Tasa de variación de la población en el nivel superior por decaimiento natural.
						\item $A_{21}$: Coeficiente de probabilidad de emisión espontánea de Einstein.
						\item $n_2$: Número de átomos en el estado excitado.
					\end{itemize}
			\end{itemize}

	\section{Teoría de circuitos}
		\subsection*{Fórmulas básicas}
			\begin{itemize}
				\item \textbf{(2.1)} $\boxed{R=\rho\frac{L}{S}}$ Resistencia eléctrica de un conductor en función de sus características físicas.
					\begin{itemize}
						\item $R$: Resistencia eléctrica. Oposición que presenta el material al paso de la corriente. Se mide en Ohmios ($\Omega$).
						\item $\rho$: Resistividad del material. Propiedad intrínseca del material del conductor. Se mide en Ohmios-metro ($\Omega\cdot m$).
						\item $L$: Longitud del conductor. Se mide en metros ($m$).
						\item $S$: Sección transversal o área del conductor. Se mide en metros cuadrados ($m^2$).
					\end{itemize}    

				\item \textbf{(2.2)} $\boxed{G=\frac{1}{R}}$ Conductancia eléctrica. Es la facilidad que ofrece un material al paso de la corriente eléctrica (el inverso de la resistencia).
					\begin{itemize}
						\item $G$: Conductancia. Se mide en Siemens ($S$) o mhos ($\Omega^{-1}$).
						\item $R$: Resistencia eléctrica. Se mide en Ohmios ($\Omega$).
					\end{itemize}

				\item \textbf{(2.3)} $\boxed{V=I\cdot R}$ Ley de Ohm. Relaciona la tensión, la corriente y la resistencia en un circuito.
					\begin{itemize}
						\item $V$: Tensión, diferencia de potencial o voltaje. Se mide en Voltios ($V$).
						\item $I$: Intensidad de corriente eléctrica. Se mide en Amperios ($A$).
						\item $R$: Resistencia eléctrica. Se mide en Ohmios ($\Omega$).
					\end{itemize}
			\end{itemize}

		\subsection*{Condensador}
			\begin{itemize}
				\item \textbf{(2.4)} $\boxed{C=\frac{Q}{V}}$ Definición de capacidad (o capacitancia). Relación entre la carga almacenada y la tensión.
					\begin{itemize}
						\item $C$: Capacidad del condensador. Se mide en Faradios ($F$).
						\item $Q$: Carga eléctrica almacenada. Se mide en Culombios ($C$).
						\item $V$: Tensión o diferencia de potencial entre las placas del condensador. Se mide en Voltios ($V$).
					\end{itemize}

				\item \textbf{(2.5)} $\boxed{i(t)=C\frac{dv(t)}{dt}}$ Relación tensión-corriente en un condensador. La corriente es proporcional a la velocidad de cambio de la tensión.
					\begin{itemize}
						\item $i(t)$: Intensidad de corriente instantánea en función del tiempo. Se mide en Amperios ($A$).
						\item $C$: Capacidad del condensador. Se mide en Faradios ($F$).
						\item $\frac{dv(t)}{dt}$: Derivada de la tensión respecto al tiempo (tasa de cambio temporal del voltaje). Se mide en Voltios por segundo ($V/s$).
					\end{itemize}

				\item \textbf{(2.6)} $\boxed{v_c(t)=\frac{1}{C}\int_{-\infty}^{t} i(t)dt}$ Tensión instantánea en un condensador a partir de la corriente.
					\begin{itemize}
						\item $v_c(t)$: Tensión instantánea en los bornes del condensador en función del tiempo. Se mide en Voltios ($V$).
						\item $C$: Capacidad del condensador. Se mide en Faradios ($F$).
						\item $i(t)$: Intensidad de corriente instantánea. Se mide en Amperios ($A$).
						\item $t$: Instante de tiempo evaluado. Se mide en segundos ($s$).
					\end{itemize}

				\item \textbf{(2.7)} $\boxed{i(t)=\frac{V}{R}e^{-\frac{1}{RC}t}}$ Corriente transitoria en un circuito RC (aplicable típicamente a la carga o descarga).
					\begin{itemize}
						\item $i(t)$: Intensidad de corriente instantánea. Se mide en Amperios ($A$).
						\item $V$: Tensión de la fuente o tensión inicial. Se mide en Voltios ($V$).
						\item $R$: Resistencia del circuito. Se mide en Ohmios ($\Omega$).
						\item $C$: Capacidad del condensador. Se mide en Faradios ($F$).
						\item $t$: Tiempo transcurrido desde el inicio del evento. Se mide en segundos ($s$).
						\item $e$: Número de Euler (base de los logaritmos neperianos, $\approx 2.718$).
					\end{itemize}

				\item \textbf{(2.8)} $\boxed{v_c(t)=V(1-e^{-\frac{1}{RC}t})}$ Ecuación de la tensión durante la \textbf{carga} de un condensador en un circuito RC.
					\begin{itemize}
						\item $v_c(t)$: Tensión instantánea en el condensador. Se mide en Voltios ($V$).
						\item $V$: Tensión de la fuente de alimentación. Se mide en Voltios ($V$).
						\item $R$: Resistencia del circuito. Se mide en Ohmios ($\Omega$).
						\item $C$: Capacidad del condensador. Se mide en Faradios ($F$).
						\item $t$: Tiempo transcurrido. Se mide en segundos ($s$).
					\end{itemize}

				\item \textbf{(2.9)} $\boxed{\tau=RC}$ Constante de tiempo de un circuito RC. Indica la rapidez con la que el condensador se carga o descarga.
					\begin{itemize}
						\item $\tau$ (tau): Constante de tiempo. Representa el tiempo necesario para que la magnitud alcance el $\approx 63.2\%$ de su valor final. Se mide en segundos ($s$).
						\item $R$: Resistencia del circuito. Se mide en Ohmios ($\Omega$).
						\item $C$: Capacidad del condensador. Se mide en Faradios ($F$).
					\end{itemize}

				\item \textbf{(2.10)} $\boxed{v_c(t)=Ve^{-\frac{1}{RC}t}}$ Ecuación de la tensión durante la \textbf{descarga} de un condensador en un circuito RC.
					\begin{itemize}
						\item $v_c(t)$: Tensión instantánea en el condensador. Se mide en Voltios ($V$).
						\item $V$: Tensión inicial almacenada en el condensador antes de la descarga. Se mide en Voltios ($V$).
						\item $R$: Resistencia del circuito. Se mide en Ohmios ($\Omega$).
						\item $C$: Capacidad del condensador. Se mide en Faradios ($F$).
						\item $t$: Tiempo transcurrido desde el inicio de la descarga. Se mide en segundos ($s$).
					\end{itemize}
			\end{itemize}

		\subsection*{Bobina}
			\begin{itemize}
				\item \textbf{(2.11)} $\boxed{v_L(t)=L\frac{di(t)}{dt}}$ Relación tensión-corriente en una bobina. La tensión es proporcional a la velocidad de cambio de la corriente.
					\begin{itemize}
						\item $v_L(t)$: Tensión instantánea en los bornes de la bobina. Se mide en Voltios ($V$).
						\item $L$: Inductancia de la bobina. \textit{(Nota: No confundir con la longitud de la fórmula 2.1)}. Se mide en Henrios ($H$).
						\item $\frac{di(t)}{dt}$: Derivada de la corriente respecto al tiempo (tasa de cambio temporal de la corriente). Se mide en Amperios por segundo ($A/s$).
					\end{itemize}

				\item \textbf{(2.12)} $\boxed{i(t)=\frac{1}{L}\int_{-\infty}^{t} v(t)dt}$ Corriente instantánea en una bobina a partir de la tensión.
					\begin{itemize}
						\item $i(t)$: Intensidad de corriente instantánea. Se mide en Amperios ($A$).
						\item $L$: Inductancia de la bobina. Se mide en Henrios ($H$).
						\item $v(t)$: Tensión instantánea en la bobina. Se mide en Voltios ($V$).
						\item $t$: Instante de tiempo evaluado. Se mide en segundos ($s$).
					\end{itemize}

				\item \textbf{(2.13)} $\boxed{i(t)=\frac{V}{R}(1-e^{-\frac{R}{L}t})}$ Ecuación de la corriente durante la \textbf{conexión} (establecimiento de corriente) en un circuito RL.
					\begin{itemize}
						\item $i(t)$: Intensidad de corriente instantánea. Se mide en Amperios ($A$).
						\item $V$: Tensión de la fuente de alimentación. Se mide en Voltios ($V$).
						\item $R$: Resistencia del circuito. Se mide en Ohmios ($\Omega$).
						\item $L$: Inductancia de la bobina. Se mide en Henrios ($H$).
						\item $t$: Tiempo transcurrido. Se mide en segundos ($s$).
					\end{itemize}

				\item \textbf{(2.14)} $\boxed{v_L(t)=Ve^{-\frac{R}{L}t}}$ Tensión transitoria en la bobina durante el establecimiento o corte de corriente.
					\begin{itemize}
						\item $v_L(t)$: Tensión instantánea en la bobina. Se mide en Voltios ($V$).
						\item $V$: Tensión inicial en el momento de la conmutación. Se mide en Voltios ($V$).
						\item $R$: Resistencia del circuito. Se mide en Ohmios ($\Omega$).
						\item $L$: Inductancia de la bobina. Se mide en Henrios ($H$).
						\item $t$: Tiempo transcurrido. Se mide en segundos ($s$).
					\end{itemize}

				\item \textbf{(2.15)} $\boxed{\tau=\frac{L}{R}}$ Constante de tiempo de un circuito RL. Indica la rapidez con la que se establece o decae la corriente en la bobina.
				\begin{itemize}
					\item $\tau$ (tau): Constante de tiempo. Se mide en segundos ($s$).
					\item $L$: Inductancia de la bobina. Se mide en Henrios ($H$).
					\item $R$: Resistencia del circuito. Se mide en Ohmios ($\Omega$).
				\end{itemize}

				\item \textbf{(2.16)} $\boxed{i(t)=\frac{V}{R}e^{-\frac{R}{L}t}}$ Ecuación de la corriente durante la \textbf{desconexión} (decaimiento de corriente) en un circuito RL.
				\begin{itemize}
					\item $i(t)$: Intensidad de corriente instantánea en decaimiento. Se mide en Amperios ($A$).
					\item $V$: Tensión inicial (donde $\frac{V}{R}$ representa la corriente máxima inicial antes del corte). Se mide en Voltios ($V$).
					\item $R$: Resistencia del circuito. Se mide en Ohmios ($\Omega$).
					\item $L$: Inductancia de la bobina. Se mide en Henrios ($H$).
					\item $t$: Tiempo transcurrido desde el corte. Se mide en segundos ($s$).
				\end{itemize}
			\end{itemize}

	\section{Electrostática}
		\subsection*{Fórmulas}
			\begin{itemize}
				\item \textbf{(3.1)} $\boxed{\vec{E}=\frac{1}{4\pi\epsilon_{0}}q^{\prime}\frac{(\vec{r}-\vec{r^{\prime}})}{||\vec{r}-\vec{r}^{\prime}||^{3}} = \frac{1}{4\pi\epsilon_{0}}q^{\prime}\frac{1}{||\vec{r}-\vec{r}^{\prime}||^{2}}\hat{u}_{\vec{r}-\vec{r'}}}$ Campo eléctrico generado por una carga puntual $q'$ en un punto del espacio.
					\begin{itemize}
						\item $\vec{E}$: Vector campo eléctrico. Se mide en Voltios por metro ($V/m$) o Newtons por Culombio ($N/C$).
						\item $\epsilon_{0}$: Permitividad eléctrica del vacío. Constante física cuyo valor es $8,854\cdot10^{-12}\,\text{F/m}$. Describe la resistencia que ofrece el vacío a la formación de un campo eléctrico.
						\item $q'$: Carga eléctrica puntual que genera el campo. Se mide en Culombios (C).
						\item $r$: Vector posición del punto donde se evalúa el campo eléctrico. Se mide en metros (m).
						\item $r'$: Vector posición de la carga generadora $q'$. Se mide en metros (m).
						\item $||r -  r'||$: Distancia en línea recta entre la carga generadora y el punto de evaluación. Se mide en metros (m).
						\item $\hat{u}_{r -  r'}$: Vector unitario que apunta desde la carga $q'$ hacia el punto de evaluación $r$. Es adimensional.
					\end{itemize}

				\item \textbf{(3.2)} $\boxed{\vec{E} = \int \mathrm{d}\vec{E} = \int \frac{1}{4\pi\epsilon_{0}} \mathrm{d}q' \frac{(\vec{r} - \vec{r}')}{\vert{}\vert{}\vec{r} - \vec{r}'\vert{}\vert{}^3}}$ Campo eléctrico generado por una distribución continua de carga (expresado con el vector diferencia de posiciones).
					\begin{itemize}
						\item $\mathrm{d}q'$: Diferencial de carga (una porción infinitesimal de la distribución de carga). Se mide en Culombios ($C$).
						\item $\mathrm{d} \vec{E}$: Diferencial de campo eléctrico. Cada una de esas cargas pequeñísimas $\mathrm{d}q'$ genera su propio "mini" campo eléctrico en el espacio. Dicho campo diminuto es $\mathrm{d}\vec{E}$. Se mide en Voltios por metro ($V/m$) o Newtons por Culombio ($N/C$).
						\item El resto es igual que en la fórmula (3.1).
					\end{itemize}

				\item \textbf{(3.3)} $\boxed{\vec{E} = \int \mathrm{d}\vec{E} = \int \frac{1}{4\pi\epsilon_{0}} \mathrm{d}q' \frac{\hat{u}_{\vec{r}-\vec{r}'}}{\vert{}\vert{}\vec{r} - \vec{r}'\vert{}\vert{}^2}}$ Campo eléctrico generado por una distribución continua de carga (expresado con el vector unitario direccional).
					\begin{itemize}
						\item Los símbolos tienen exactamente el mismo significado que en las fórmulas (3.1) y (3.2).
					\end{itemize}

				\item \textbf{(3.4)} $\boxed{\mathrm{d}q' = \lambda \mathrm{d}l}$ Relación del diferencial de carga para una distribución \textbf{lineal} (1D).
					\begin{itemize}
						\item $\mathrm{d}q'$: Diferencial de carga. Se mide en Culombios ($C$).
						\item $\lambda$: Densidad lineal de carga. Cantidad de carga por unidad de longitud. Se mide en Culombios por metro ($C/m$).
						\item $\mathrm{d}l$: Diferencial de longitud. Se mide en metros ($m$).
					\end{itemize}

				\item \textbf{(3.5)} $\boxed{\mathrm{d}q' = \sigma \mathrm{d}S}$ Relación del diferencial de carga para una distribución \textbf{superficial} (2D).
					\begin{itemize}
						\item $\mathrm{d}q'$: Diferencial de carga. Se mide en Culombios ($C$).
						\item $\sigma$: Densidad superficial de carga. Cantidad de carga por unidad de área. Se mide en Culombios por metro cuadrado ($C/m^2$).
						\item $\mathrm{d}S$: Diferencial de superficie o área. Se mide en metros cuadrados ($m^2$).
					\end{itemize}

				\item \textbf{(3.6)} $\boxed{\mathrm{d}q' = \rho \mathrm{d}V}$ Relación del diferencial de carga para una distribución \textbf{volumétrica} (3D).
					\begin{itemize}
						\item $\mathrm{d}q'$: Diferencial de carga. Se mide en Culombios ($C$).
						\item $\rho$: Densidad volumétrica de carga. Cantidad de carga por unidad de volumen. Se mide en Culombios por metro cúbico ($C/m^3$).
						\item $\mathrm{d}V$: Diferencial de volumen. Se mide en metros cúbicos ($m^3$).
					\end{itemize}

				\item \textbf{(3.7)} $\boxed{\vec{F}_e = \frac{1}{4\pi\epsilon_{0}} \frac{q q'}{d^2} \hat{u}_{\vec{r}-\vec{r}'} = \frac{1}{4\pi\epsilon_{0}} \frac{q q'}{d^2} \frac{(\vec{r} - \vec{r}')}{\vert{}\vert{}\vec{r} - \vec{r}'\vert{}\vert{}}}$ Ley de Coulomb. Fuerza electrostática entre dos cargas puntuales.
					\begin{itemize}
						\item $\vec{F}_e$: Vector fuerza electrostática. Se mide en Newtons ($N$).
						\item $\epsilon_{0}$: Permitividad eléctrica del vacío. Constante física cuyo valor es $8,854\cdot10^{-12}\,\text{F/m}$. Describe la resistencia que ofrece el vacío a la formación de un campo eléctrico.
						\item $q, q'$: Valor de las cargas eléctricas que interactúan. Se miden en Culombios ($C$).
						\item $d$: Distancia que separa a ambas cargas (equivalente a $\vert{}\vert{}\vec{r} - \vec{r}'\vert{}\vert{}$). Se mide en metros ($m$).
						\item $\hat{u}_{\vec{r}-\vec{r}'}$: Vector unitario en la dirección que une ambas cargas.
					\end{itemize}

				\item \textbf{(3.8)} $\boxed{\vec{F}_e = q \vec{E}(\vec{r})}$ Fuerza electrostática que experimenta una carga $q$ al situarse en un campo eléctrico externo $\vec{E}$.
					\begin{itemize}
						\item $\vec{F}_e$: Vector fuerza electrostática. Se mide en Newtons ($N$).
						\item $q$: Carga eléctrica puntual que interactúa con el campo electromagnético. Se mide en Culombios ($C$).
						\item $\vec{E}(\vec{r})$: Vector campo eléctrico en la posición $\vec{r}$ donde se encuentra la carga. Se mide en Voltios por metro ($V/m$) o Newtons por Culombio ($N/C$).
					\end{itemize}

				\item \textbf{(3.9)} $\boxed{\vec{F} = \int_\Gamma \mathrm{d}q' \vec{E}(\vec{r}')}$ Fuerza electrostática total sobre una distribución continua de carga inmersa en un campo eléctrico.
					\begin{itemize}
						\item $\vec{F}$: Vector fuerza neta total. Se mide en Newtons ($N$).
						\item $\mathrm{d}q'$: Diferencial de carga de la distribución. Se mide en Culombios ($C$).
						\item $\vec{E}(\vec{r}')$: Vector campo eléctrico evaluado en la posición del diferencial de carga. Se mide en Voltios por metro ($V/m$) o Newtons por Culombio ($N/C$).
					\end{itemize}

				\item \textbf{(3.10)} $\boxed{\phi_E = \int_S \vec{E} \cdot \mathrm{d}\vec{S}}$ Flujo del campo eléctrico a través de una superficie abierta o cerrada $S$.
					\begin{itemize}
						\item $\phi_E$: Flujo eléctrico. Representa la "cantidad" de líneas de campo que atraviesan una superficie. Se mide en Voltios-metro ($V\cdot m$) o Newtons-metro cuadrado por Culombio ($N\cdot m^2/C$).
						\item $\vec{E}$: Vector campo eléctrico. Se mide en Voltios por metro ($V/m$).
						\item $\mathrm{d}\vec{S}$: Vector diferencial de superficie (su dirección es normal a la superficie en ese punto). Se mide en metros cuadrados ($m^2$).
					\end{itemize}

				\item \textbf{(3.11)} $\boxed{\oint_S \vec{E} \cdot \mathrm{d}\vec{S} = \frac{Q_{int}}{\epsilon_0}}$ Ley de Gauss. El flujo eléctrico total a través de una superficie cerrada es proporcional a la carga neta encerrada.
					\begin{itemize}
						\item $\oint_S \vec{E} \cdot \mathrm{d}\vec{S}$: Flujo eléctrico a través de una superficie cerrada $S$. Se mide en Voltios-metro ($V\cdot m$).
						\item $Q_{int}$: Carga eléctrica neta encerrada dentro de la superficie cerrada $S$. Se mide en Culombios ($C$).
						\item $\epsilon_{0}$: Permitividad eléctrica del vacío. Constante física cuyo valor es $8,854\cdot10^{-12}\,\text{F/m}$. Describe la resistencia que ofrece el vacío a la formación de un campo eléctrico.
					\end{itemize}

				\item \textbf{(3.12)} $\boxed{V(\vec{r}) = \frac{1}{4\pi\epsilon_0} q' \frac{1}{\vert{}\vert{}\vec{r} - \vec{r}'\vert{}\vert{}}}$ Potencial eléctrico generado por una carga puntual $q'$ en un punto $\vec{r}$.
					\begin{itemize}
						\item $V(\vec{r})$: Potencial eléctrico (voltaje) en la posición $\vec{r}$. Se mide en Voltios ($V$).
						\item $\epsilon_{0}$: Permitividad eléctrica del vacío. Constante física cuyo valor es $8,854\cdot10^{-12}\,\text{F/m}$. Describe la resistencia que ofrece el vacío a la formación de un campo eléctrico.
						\item $q'$: Carga eléctrica puntual generadora. Se mide en Culombios ($C$).
						\item $\vert{}\vert{}\vec{r} - \vec{r}'\vert{}\vert{}$: Distancia entre la carga y el punto de evaluación. Se mide en metros ($m$).
					\end{itemize}

				\item \textbf{(3.13)} $\boxed{V(\vec{r}) = \int_\Omega \frac{1}{4\pi\epsilon_0} \mathrm{d}q' \frac{1}{\vert{}\vert{}\vec{r} - \vec{r}'\vert{}\vert{}}}$ Potencial eléctrico generado por una distribución continua de carga.
					\begin{itemize}
						\item $V(\vec{r})$: Potencial eléctrico total. Se mide en Voltios ($V$).
						\item $\mathrm{d}q'$: Diferencial de carga de la distribución. Se mide en Culombios ($C$).
						\item $\vert{}\vert{}\vec{r} - \vec{r}'\vert{}\vert{}$: Distancia entre el diferencial de carga y el punto evaluado. Se mide en metros ($m$).
					\end{itemize}

				\item \textbf{(3.14)} $\boxed{U = \frac{1}{4\pi\epsilon_0} \frac{1}{2} \sum_{i,j=1 (i\neq j)}^N \frac{q_i q_j}{d_{ij}}}$ Energía potencial electrostática almacenada en un sistema de $N$ cargas puntuales.
					\begin{itemize}
						\item $U$: Energía potencial electrostática. Se mide en Julios ($J$).
						\item $\epsilon_{0}$: Permitividad eléctrica del vacío. Constante física cuyo valor es $8,854\cdot10^{-12}\,\text{F/m}$. Describe la resistencia que ofrece el vacío a la formación de un campo eléctrico.
						\item $q_i, q_j$: Valor de los distintos pares de cargas del sistema. Se miden en Culombios ($C$).
						\item $d_{ij}$: Distancia de separación entre la carga $i$ y la carga $j$. Se mide en metros ($m$).
					\end{itemize}

				\item \textbf{(3.15)} $\boxed{U = \frac{1}{2} \int_\Omega \mathrm{d}q' V(\vec{r}') = \frac{\epsilon_0}{2} \int_T \vert{}\vert{}\vec{E}(\vec{r})\vert{}\vert{}^2 \mathrm{d}^3\vec{r}}$ Relaciones de la energía potencial electrostática para distribuciones continuas y energía almacenada en el propio campo eléctrico.
					\begin{itemize}
						\item $U$: Energía potencial electrostática total. Se mide en Julios ($J$).
						\item $\mathrm{d}q'$: Diferencial de carga. Se mide en Culombios ($C$).
						\item $V(\vec{r}')$: Potencial eléctrico en la posición del diferencial de carga. Se mide en Voltios ($V$).
						\item $\epsilon_{0}$: Permitividad eléctrica del vacío. Constante física cuyo valor es $8,854\cdot10^{-12}\,\text{F/m}$. Describe la resistencia que ofrece el vacío a la formación de un campo eléctrico.
						\item $\vert{}\vert{}\vec{E}(\vec{r})\vert{}\vert{}$: Módulo del campo eléctrico en cada punto del espacio. Se mide en Voltios por metro ($V/m$).
						\item $\mathrm{d}^3\vec{r}$: Elemento diferencial de volumen en el espacio tridimensional (equivalente a $\mathrm{d}V$). Se mide en metros cúbicos ($m^3$).
					\end{itemize}

				\item \textbf{(3.16)} $\boxed{E = \frac{\sigma}{\epsilon_0}}$ Magnitud del campo eléctrico en las proximidades de la superficie de un conductor ideal en equilibrio o entre las placas de un condensador plano.
					\begin{itemize}
						\item $E$: Módulo del campo eléctrico. Se mide en Voltios por metro ($V/m$) o Newtons por Culombio ($N/C$).
						\item $\sigma$: Densidad superficial de carga. Se mide en Culombios por metro cuadrado ($C/m^2$).
						\item $\epsilon_{0}$: Permitividad eléctrica del vacío. Constante física cuyo valor es $8,854\cdot10^{-12}\,\text{F/m}$. Describe la resistencia que ofrece el vacío a la formación de un campo eléctrico.
					\end{itemize}

				\item \textbf{(3.17)} $\boxed{\vert{}\Delta V\vert{} = \frac{\sigma}{\epsilon_0} d = \frac{Q}{\epsilon_0 S} d}$ Diferencia de potencial (voltaje) entre dos placas plano-paralelas.
					\begin{itemize}
						\item $\vert{}\Delta V\vert{}$: Valor absoluto de la diferencia de potencial o caída de tensión. Se mide en Voltios ($V$).
						\item $\sigma$: Densidad superficial de carga en las placas. Se mide en Culombios por metro cuadrado ($C/m^2$).
						\item $\epsilon_{0}$: Permitividad eléctrica del vacío. Constante física cuyo valor es $8,854\cdot10^{-12}\,\text{F/m}$. Describe la resistencia que ofrece el vacío a la formación de un campo eléctrico.
						\item $d$: Distancia de separación entre las placas. Se mide en metros ($m$).
						\item $Q$: Carga total almacenada en una de las placas. Se mide en Culombios ($C$).
						\item $S$: Área de la superficie de las placas. Se mide en metros cuadrados ($m^2$).
					\end{itemize}

				\item \textbf{(3.18)} $\boxed{C = \frac{Q}{\vert{}\Delta V\vert{}} = \frac{\epsilon_0}{d} S}$ Capacidad de un condensador de placas plano-paralelas ideal.
					\begin{itemize}
						\item $C$: Capacidad eléctrica (capacitancia). Se mide en Faradios ($F$).
						\item $Q$: Carga eléctrica almacenada en el condensador. Se mide en Culombios ($C$).
						\item $\vert{}\Delta V\vert{}$: Diferencia de potencial entre las placas. Se mide en Voltios ($V$).
						\item $\epsilon_{0}$: Permitividad eléctrica del vacío. Constante física cuyo valor es $8,854\cdot10^{-12}\,\text{F/m}$. Describe la resistencia que ofrece el vacío a la formación de un campo eléctrico.
						\item $d$: Distancia de separación entre las placas. Se mide en metros ($m$).
						\item $S$: Área de la superficie de una placa. Se mide en metros cuadrados ($m^2$).
					\end{itemize}
			\end{itemize}

	\section{Magnetostática e Inducción electromagnética}
		\subsection*{Fórmulas}
			\begin{itemize}
				\item \textbf{(4.1)} $\boxed{\vec{B}=\frac{\mu_0}{4 \pi} \frac{q \vec{v} \times (\vec{r} - \vec{r'})}{||\vec{r}-\vec{r'}||^3}=\frac{\mu_{0}}{4\pi}\frac{q\vec{v}\times\hat{u}_{\vec{r}-\vec{r'}}}{||\vec{r}-\vec{r^{\prime}}||^{2}}}$ Ley de Biot-Savart para una carga puntual en movimiento. Campo magnético generado por una carga en un punto del espacio.
					\begin{itemize}
						\item $\vec{B}$: Vector campo magnético. Se mide en Teslas ($T$).
						\item $\mu_0$: Permeabilidad magnética del vacío. Constante física cuyo valor es $4 \pi \times 10^{-7} \; T \cdot m/A$.
						\item $q$: Carga eléctrica puntual que interactúa con el campo electromagnético. Se mide en Culombios ($C$).
						\item $\vec{v}$: Vector velocidad de la carga. Se mide en metros por segundo ($m/s$).
						\item $\vec{r}$: Vector posición del punto donde se evalúa el campo magnético. Se mide en metros ($m$). 
						\item $\vec{r'}$: Vector posición de la carga generadora. Se mide en metros ($m$).
						\item $||\vec{r} -  \vec{r'}||$: Distancia en línea recta entre la carga y el punto de evaluación. Se mide en metros ($m$).
						\item $u^r -  r'$: Vector unitario que apunta desde la carga hacia el punto de evaluación. Es adimensional.
					\end{itemize}

				\item \textbf{(4.2)} $\boxed{\vec{B} = \int_{\mathcal{C}} \frac{\mu_0}{4\pi} \frac{I\mathrm{d}\vec{l} \times (\vec{r} - \vec{r}')}{\vert{}\vert{}\vec{r} - \vec{r}'\vert{}\vert{}^3} = \int_{\mathcal{C}} \frac{\mu_0}{4\pi} \frac{I\mathrm{d}\vec{l} \times \hat{u}_{\vec{r}-\vec{r}'}}{\vert{}\vert{}\vec{r} - \vec{r}'\vert{}\vert{}^2}}$ Ley de Biot-Savart para un hilo conductor. Campo magnético generado por una corriente eléctrica a lo largo de un circuito $\mathcal{C}$.
					\begin{itemize}
						\item $\vec{B}$: Vector campo magnético total. Se mide en Teslas ($T$).
						\item $\mu_0$: Permeabilidad magnética del vacío. Constante física cuyo valor es $4 \pi \times 10^{-7} \; T \cdot m/A$.
						\item $I$: Intensidad de corriente eléctrica que circula por el sistema. Se mide en Amperios ($A$).
						\item $\mathrm{d}\vec{l}$: Vector diferencial de longitud del circuito, tangente al mismo en el sentido de la corriente. Se mide en metros ($m$).
						\item $\vec{r}, \vec{r}', \vert{}\vert{}\vec{r} - \vec{r}'\vert{}\vert{}, \hat{u}_{\vec{r}-\vec{r}'}$: Mismo significado que en la fórmula (4.1).
					\end{itemize}

				\item \textbf{(4.3)} $\boxed{\oint_{\Gamma} \vec{B} \cdot \mathrm{d}\vec{l} = \mu_0 I_{int}}$ Ley de Ampère. La circulación del campo magnético a lo largo de una curva cerrada $\Gamma$ es proporcional a la corriente neta que la atraviesa.
					\begin{itemize}
						\item $\oint_{\Gamma} \vec{B} \cdot \mathrm{d}\vec{l}$: Circulación del campo magnético a lo largo de una trayectoria cerrada $\Gamma$. Se mide en Teslas-metro ($T\cdot m$).
						\item $\vec{B}$: Vector campo magnético. Se mide en Teslas ($T$).
						\item $\mathrm{d}\vec{l}$: Vector diferencial de camino a lo largo de la curva $\Gamma$. Se mide en metros ($m$).
						\item $\mu_0$: Permeabilidad magnética del vacío.
						\item $I$: Intensidad de corriente eléctrica neta encerrada que circula por el sistema. Se mide en Amperios ($A$).
					\end{itemize}

				\item \textbf{(4.4)} $\boxed{\vec{F}_{1\to2} = \frac{\mu_0}{4\pi} \frac{q_1 q_2}{d^2} (\vec{v}_2 \times (\vec{v}_1 \times \hat{u}_{\vec{r}_2-\vec{r}_1}))}$ Fuerza magnética que ejerce una carga en movimiento $q_1$ sobre otra carga en movimiento $q_2$.
					\begin{itemize}
						\item $\vec{F}_{1\to2}$: Vector fuerza magnética ejercida por la carga 1 sobre la carga 2. Se mide en Newtons ($N$).
						\item $\mu_0$: Permeabilidad magnética del vacío.
						\item $q_1, q_2$: Valor de las cargas eléctricas que interactúan. Se miden en Culombios ($C$).
						\item $d$: Distancia que separa a ambas cargas. Se mide en metros ($m$).
						\item $\vec{v}_1, \vec{v}_2$: Vectores velocidad de la carga 1 y 2, respectivamente. Se miden en metros por segundo ($m/s$).
						\item $\hat{u}_{\vec{r}_2-\vec{r}_1}$: Vector unitario que apunta desde la carga 1 hacia la carga 2. Es adimensional.
					\end{itemize}

				\item \textbf{(4.5)} $\boxed{\vec{F}_m = q\vec{v} \times \vec{B}(\vec{r})}$ Fuerza magnética (componente magnética de la fuerza de Lorentz) sobre una carga puntual que se mueve en un campo magnético externo.
					\begin{itemize}
						\item $\vec{F}_m$: Vector fuerza magnética. Se mide en Newtons ($N$).
						\item $q$: Carga eléctrica puntual que interactúa con el campo electromagnético. Se mide en Culombios ($C$).
						\item $\vec{v}$: Vector velocidad de la carga. Se mide en metros por segundo ($m/s$).
						\item $\vec{B}(\vec{r})$: Vector campo magnético en la posición de la carga. Se mide en Teslas ($T$).
					\end{itemize}

				\item \textbf{(4.6)} $\boxed{\vec{F}_m = \int_{\Gamma} I\mathrm{d}\vec{l} \times \vec{B}(\vec{r})}$ Fuerza de Laplace. Fuerza magnética total ejercida sobre un hilo conductor recorrido por una corriente e inmerso en un campo magnético.
					\begin{itemize}
						\item $\vec{F}_m$: Vector fuerza magnética neta. Se mide en Newtons ($N$).
						\item $I$: Intensidad de corriente eléctrica que circula por el sistema. Se mide en Amperios ($A$).
						\item $\mathrm{d}\vec{l}$: Vector diferencial de longitud del hilo. Se mide en metros ($m$).
						\item $\vec{B}(\vec{r})$: Vector campo magnético externo evaluado en cada punto del hilo. Se mide en Teslas ($T$).
					\end{itemize}

				\item \textbf{(4.7)} $\boxed{\vec{F}_{em} = q\left[\vec{E}(\vec{r}) + \vec{v} \times \vec{B}(\vec{r})\right]}$ Fuerza de Lorentz completa. Es la fuerza electromagnética total sobre una carga puntual en presencia de un campo eléctrico y magnético simultáneamente.
					\begin{itemize}
						\item $\vec{F}_{em}$: Vector fuerza electromagnética total. Se mide en Newtons ($N$).
						\item $q$: Carga eléctrica puntual que interactúa con el campo electromagnético. Se mide en Culombios ($C$).
						\item $\vec{E}(\vec{r})$: Vector campo eléctrico en la posición de la carga. Se mide en Voltios por metro ($V/m$) o Newtons por Culombio ($N/C$).
						\item $\vec{v}$: Vector velocidad de la carga. Se mide en metros por segundo ($m/s$).
						\item $\vec{B}(\vec{r})$: Vector campo magnético en la posición de la carga. Se mide en Teslas ($T$).
					\end{itemize}

				\item \textbf{(4.8)} $\boxed{\phi_M = \int_S \vec{B} \cdot \mathrm{d}\vec{S}}$ Flujo magnético a través de una superficie $S$.
					\begin{itemize}
						\item $\phi_M$: Flujo magnético. Representa la cantidad de líneas de campo magnético que atraviesan una superficie. Se mide en Webers ($Wb$) o Teslas por metro cuadrado ($T\cdot m^2$).
						\item $\vec{B}$: Vector campo magnético. Se mide en Teslas ($T$).
						\item $\mathrm{d}\vec{S}$: Vector diferencial de superficie (su dirección es perpendicular a la superficie). Se mide en metros cuadrados ($m^2$).
					\end{itemize}

				\item \textbf{(4.9)} $\boxed{\text{f.e.m} = -\frac{\mathrm{d}\phi_B}{\mathrm{d}t}}$ Ley de Faraday-Lenz. La fuerza electromotriz inducida en un circuito cerrado es directamente proporcional a la rapidez con la que cambia el flujo magnético que lo atraviesa (el signo negativo indica oposición al cambio).
					\begin{itemize}
						\item $\text{f.e.m}$: Fuerza electromotriz inducida (voltaje inducido). Se mide en Voltios ($V$).
						\item $\mathrm{d}\phi_B$: Diferencial de flujo magnético. Se mide en Webers ($Wb$). \textit{(Nota: $\phi_B$ es lo mismo que $\phi_M$ en la fórmula anterior)}.
						\item $\mathrm{d}t$: Diferencial de tiempo. Se mide en segundos ($s$).
					\end{itemize}

				\item \textbf{(4.10)} $\boxed{B = \mu_0 n I}$ Magnitud del campo magnético en el interior de un solenoide ideal.
					\begin{itemize}
						\item $B$: Módulo del campo magnético. Se mide en Teslas ($T$).
						\item $\mu_0$: Permeabilidad magnética del vacío.
						\item $n$: Densidad de espiras del solenoide (número de vueltas dividido por su longitud, $n = N/L$). Se mide en metros a la menos uno ($m^{-1}$).
						\item $I$: Intensidad de corriente eléctrica que circula por el sistema. Se mide en Amperios ($A$).
					\end{itemize}

				\item \textbf{(4.11)} $\boxed{L = \frac{d\Phi}{dI}}$ Definición de inductancia o coeficiente de autoinducción de una bobina/circuito.
					\begin{itemize}
						\item $L$: Inductancia o coeficiente de autoinducción. Se mide en Henrios ($H$).
						\item $d\Phi$: Cambio en el flujo magnético propio. Se mide en Webers ($Wb$). \textit{(Nota: $\Phi$ es la misma magnitud que $\phi_M$ o $\phi_B$)}.
						\item $dI$: Cambio en la intensidad de corriente. Se mide en Amperios ($A$).
					\end{itemize}

				\item \textbf{(4.12)} $\boxed{\varepsilon = -L \frac{dI}{dt}}$ Fuerza electromotriz autoinducida en una bobina debido a un cambio en la corriente que la atraviesa.
					\begin{itemize}
						\item $\varepsilon$: Fuerza electromotriz (f.e.m.) o tensión inducida. Se mide en Voltios ($V$). \textit{(Nota: $\varepsilon$ es equivalente a la f.e.m. de la fórmula 4.9)}.
						\item $L$: Inductancia o coeficiente de autoinducción. Se mide en Henrios ($H$).
						\item $\frac{dI}{dt}$: Tasa de cambio de la corriente eléctrica respecto al tiempo. Se mide en Amperios por segundo ($A/s$).
					\end{itemize}
			\end{itemize}

	\section{Semiconductores}
		\subsection*{Los materiales semiconductores}
			\begin{itemize}
				\item \textbf{(5.1)} $\boxed{E_{n}=\frac{-13.6}{n^{2}}\,\text{eV}}$ Energía de los niveles electrónicos en el modelo atómico de Bohr (normalmente para el átomo de hidrógeno, el cual se usa como referencia inicial).
					\begin{itemize}
						\item $E_{n}$: Energía del nivel cuántico $n$. Se mide en electronvoltios (eV).
						\item $n$: Número cuántico principal (nivel de energía). Es adimensional y su valor es un número entero.
					\end{itemize}

				\item \textbf{(5.2)} $\boxed{n_i = N_C \cdot e^{-\frac{E_C - E_F}{k_B T}}}$ Concentración intrínseca de electrones en la banda de conducción.
					\begin{itemize}
						\item $n_i$: Concentración de electrones en el material intrínseco. Se mide en electrones por centímetro cúbico ($cm^{-3}$) o por metro cúbico ($m^{-3}$).
						\item $N_C$: Densidad efectiva de estados en la banda de conducción. Indica cuántos "huecos disponibles" hay para electrones. Se mide en $cm^{-3}$ o $m^{-3}$.
						\item $E_C$: Energía del borde inferior de la banda de conducción. Se mide en Julios (J) o electronvoltios (eV).
						\item $E_F$: Energía del nivel de Fermi. Se mide en Julios (J) o electronvoltios (eV).
						\item $k_B$: Constante de Boltzmann. Su valor es aproximadamente $1.38 \times 10^{-23}$ J/K o $8.617 \times 10^{-5}$ eV/K.
						\item $T$: Temperatura absoluta del material. Se mide en Kelvin (K).
					\end{itemize}

				\item \textbf{(5.3)} $\boxed{p_i = N_V \cdot e^{-\frac{E_F - E_V}{k_B T}}}$ Concentración intrínseca de huecos en la banda de valencia.
					\begin{itemize}
						\item $p_i$: Concentración de huecos en el material intrínseco. Se mide en $cm^{-3}$ o $m^{-3}$.
						\item $N_V$: Densidad efectiva de estados en la banda de valencia. Se mide en $cm^{-3}$ o $m^{-3}$.
						\item $E_F$: Energía del nivel de Fermi. Se mide en Julios (J) o electronvoltios (eV).
						\item $E_V$: Energía del borde superior de la banda de valencia. Se mide en Julios (J) o electronvoltios (eV).
						\item $k_B$: Constante de Boltzmann. Su valor es aproximadamente $1.38 \times 10^{-23}$ J/K o $8.617 \times 10^{-5}$ eV/K.
						\item $T$: Temperatura absoluta. Se mide en Kelvin (K).
					\end{itemize}

				\item \textbf{(5.4)} $\boxed{E_F = \frac{E_C + E_V}{2} + \frac{3k_B T}{4} \ln \left( \frac{m_p}{m_n} \right)}$ Nivel de Fermi intrínseco. Muestra que el nivel de Fermi está aproximadamente en la mitad del gap de energía, con una pequeña desviación dependiente de la temperatura y las masas efectivas.
					\begin{itemize}
						\item $E_F$: Energía del nivel de Fermi intrínseco. Se mide en Julios (J) o electronvoltios (eV).
						\item $E_C$: Energía de la banda de conducción.
						\item $E_V$: Energía de la banda de valencia.
						\item $k_B$: Constante de Boltzmann. Su valor es aproximadamente $1.38 \times 10^{-23}$ J/K o $8.617 \times 10^{-5}$ eV/K.
						\item $T$: Temperatura absoluta en Kelvin (K).
						\item $m_p$: Masa efectiva de los huecos. Se mide en kilogramos (kg).
						\item $m_n$: Masa efectiva de los electrones. Se mide en kilogramos (kg).
					\end{itemize}

				\item \textbf{(5.5)} $\boxed{n \cdot p = n_i^2}$ Ley de acción de masas para semiconductores en equilibrio térmico.
					\begin{itemize}
						\item $n$: Concentración total de electrones en la banda de conducción. Se mide en $cm^{-3}$ o $m^{-3}$.
						\item $p$: Concentración total de huecos en la banda de valencia. Se mide en $cm^{-3}$ o $m^{-3}$.
						\item $n_i$: Concentración de portadores intrínsecos. Se mide en $cm^{-3}$ o $m^{-3}$.
					\end{itemize}

				\item \textbf{(5.6)} $\boxed{E_{F,n} = E_{F,i} + k_B T \cdot \ln \left( \frac{N_D}{n_i} \right)}$ Posición del nivel de Fermi en un semiconductor dopado tipo n.
					\begin{itemize}
						\item $E_{F,n}$: Energía del nivel de Fermi en el material tipo n. Se mide en eV o Julios.
						\item $E_{F,i}$: Energía del nivel de Fermi intrínseco (sin dopar). Se mide en eV o Julios.
						\item $k_B$: Constante de Boltzmann. Su valor es aproximadamente $1.38 \times 10^{-23}$ J/K o $8.617 \times 10^{-5}$ eV/K.
						\item $T$: Temperatura absoluta en Kelvin (K).
						\item $N_D$: Concentración de átomos donadores (impurezas que aportan electrones). Se mide en $cm^{-3}$ o $m^{-3}$.
						\item $n_i$: Concentración intrínseca.
					\end{itemize}

				\item \textbf{(5.7)} $\boxed{n_n = N_D, \quad p_n = \frac{n_i^2}{N_D}}$ Concentración de portadores mayoritarios y minoritarios en un semiconductor tipo n.
					\begin{itemize}
						\item $n_n$: Concentración de electrones (mayoritarios). Aproximadamente igual a la de átomos donadores $N_D$. Se mide en $cm^{-3}$.
						\item $p_n$: Concentración de huecos (minoritarios). Calculada mediante la ley de acción de masas. Se mide en $cm^{-3}$.
						\item $N_D$: Concentración de átomos donadores.
						\item $n_i$: Concentración intrínseca.
					\end{itemize}

				\item \textbf{(5.8)} $\boxed{E_{F,p} = E_{F,i} - k_B T \cdot \ln \left( \frac{N_A}{n_i} \right)}$ Posición del nivel de Fermi en un semiconductor dopado tipo p.
					\begin{itemize}
						\item $E_{F,p}$: Energía del nivel de Fermi en el material tipo p. Se mide en eV o Julios.
						\item $E_{F,i}$: Energía del nivel de Fermi intrínseco.
						\item $k_B$: Constante de Boltzmann. Su valor es aproximadamente $1.38 \times 10^{-23}$ J/K o $8.617 \times 10^{-5}$ eV/K.
						\item $T$: Temperatura absoluta en Kelvin (K).
						\item $N_A$: Concentración de átomos aceptores (impurezas que aportan huecos). Se mide en $cm^{-3}$ o $m^{-3}$.
						\item $n_i$: Concentración intrínseca.
					\end{itemize}

				\item \textbf{(5.9)} $\boxed{p_p = N_A, \quad n_p = \frac{n_i^2}{N_A}}$ Concentración de portadores mayoritarios y minoritarios en un semiconductor tipo p.
					\begin{itemize}
						\item $p_p$: Concentración de huecos (mayoritarios). Aproximadamente igual a $N_A$. Se mide en $cm^{-3}$.
						\item $n_p$: Concentración de electrones (minoritarios). Se mide en $cm^{-3}$.
						\item $N_A$: Concentración de átomos aceptores.
						\item $n_i$: Concentración intrínseca.
					\end{itemize}

				\item \textbf{(5.10)} $\boxed{\vec{v}_{d,n} = -\mu_n \cdot \vec{E}}$ Velocidad de arrastre (drift) de los electrones bajo un campo eléctrico. El signo negativo indica que se mueven en sentido contrario al campo.
					\begin{itemize}
						\item $\vec{v}_{d,n}$: Vector velocidad de arrastre de los electrones. Se mide en metros por segundo (m/s).
						\item $\mu_n$: Movilidad de los electrones. Mide la facilidad con la que se mueven por la red cristalina. Se mide en $cm^2/(V\cdot s)$ o $m^2/(V\cdot s)$.
						\item $\vec{E}$: Vector campo eléctrico aplicado. Se mide en Voltios por metro (V/m).
					\end{itemize}

				\item \textbf{(5.11)} $\boxed{\vec{v}_{d,p} = \mu_p \cdot \vec{E}}$ Velocidad de arrastre de los huecos. Se mueven en el mismo sentido que el campo eléctrico.
					\begin{itemize}
						\item $\vec{v}_{d,p}$: Vector velocidad de arrastre de los huecos. Se mide en metros por segundo (m/s).
						\item $\mu_p$: Movilidad de los huecos. Se mide en $cm^2/(V\cdot s)$ o $m^2/(V\cdot s)$.
						\item $\vec{E}$: Vector campo eléctrico aplicado. Se mide en Voltios por metro (V/m).
					\end{itemize}

				\item \textbf{(5.12)} $\boxed{I = \frac{V}{\left[ \frac{1}{q(\mu_p \cdot p + \mu_n \cdot n)} \right] \frac{L}{A}} = \frac{V}{R}}$ Ley de Ohm macroscópica para un bloque semiconductor, expresada en función de las propiedades de los portadores de carga.
					\begin{itemize}
						\item $I$: Intensidad de corriente eléctrica que circula por el sistema. Se mide en Amperios ($A$).
						\item $V$: Diferencia de potencial (voltaje) aplicada. Se mide en Voltios (V).
						\item $q$: Carga elemental del electrón ($1,602 \times 10^{-19}$ C). Se mide en Culombios (C).
						\item $\mu_p, \mu_n$: Movilidad de huecos y electrones.
						\item $p, n$: Concentración de huecos y electrones.
						\item $L$: Longitud del bloque semiconductor. Se mide en metros (m).
						\item $A$: Área de la sección transversal. Se mide en metros cuadrados ($m^2$).
						\item $R$: Resistencia total del bloque. Se mide en Ohmios ($\Omega$). El término en el corchete es la resistividad $\rho$.
					\end{itemize}

				\item \textbf{(5.13)} $\boxed{I = \frac{\Delta Q}{\tau_c} = (n_2 - n_1) \frac{q \cdot l}{2\tau_c} A = (n_2 - n_1) \frac{q \cdot v_{th}}{2} A}$ Corriente descrita mediante la transferencia neta de carga en un tiempo de tránsito, relacionada con la velocidad térmica.
					\begin{itemize}
						\item $I$: Intensidad de corriente eléctrica que circula por el sistema. Se mide en Amperios ($A$).
						\item $\Delta Q$: Variación de carga o carga neta transferida. Se mide en Culombios (C).
						\item $\tau_c$: Tiempo de colisión o de tránsito libre medio. Se mide en segundos (s).
						\item $(n_2 - n_1)$: Diferencia de concentración de portadores entre dos regiones.
						\item $q$: Carga elemental del electrón. Es una constante física cuyo valor es $1,602\cdot10^{-19}\,\text{C}$. Se mide en Culombios ($C$).
						\item $l$: Camino libre medio (distancia media entre colisiones). Se mide en metros (m).
						\item $A$: Área de la sección. Se mide en metros cuadrados ($m^2$).
						\item $v_{th}$: Velocidad térmica promedio de los portadores. Se mide en metros por segundo (m/s).
					\end{itemize}
			\end{itemize}

		\subsubsection*{La unión p-n. Los díodos}
			\begin{itemize}
				\item \textbf{(5.14)} $\boxed{V_{bi} = \frac{1}{2} \frac{q \left ( N_D x_n + N_D x_p \right )}{2 \epsilon_0 \epsilon_r} W}$ Potencial de contacto o potencial interno de la unión p-n expresado geométricamente respecto al ancho de la zona de carga espacial.
					\begin{itemize}
						\item $V_{bi}$: Potencial de contacto que se forma naturalmente en la unión. Se mide en Voltios (V).
						\item $q$: Carga elemental del electrón. Es una constante física cuyo valor es $1,602\cdot10^{-19}\,\text{C}$. Se mide en Culombios ($C$).
						\item $N_D$: Concentración de donadores (lado $n$).
						\item $N_A$: Concentración de aceptores (lado $p$).
						\item $x_n$: Extensión de la zona de vaciamiento hacia el lado n. Se mide en metros (m) o centímetros (cm).
						\item $x_p$: Extensión de la zona de vaciamiento hacia el lado p. Se mide en metros (m) o centímetros (cm).
						\item $\epsilon_{0}$: Permitividad eléctrica del vacío. Constante física cuyo valor es $8,854\cdot10^{-12}\,\text{F/m}$. Describe la resistencia que ofrece el vacío a la formación de un campo eléctrico.
						\item $\epsilon_r$: Permitividad relativa (constante dieléctrica) del semiconductor.
						\item $W$: Ancho total de la zona de vaciamiento ($W = x_n + x_p$). Se mide en metros (m).
					\end{itemize}

				\item \textbf{(5.15)} $\boxed{V_{bi} = \frac{k_B T}{q} \ln \left( \frac{N_A \cdot N_D}{n_i^2} \right)}$ Fórmula exacta fundamental para calcular el potencial de contacto $V_{bi}$ a partir de los dopajes.
					\begin{itemize}
						\item $V_{bi}$: Potencial interno. Se mide en Voltios (V).
						\item $k_B$: Constante de Boltzmann. Su valor es aproximadamente $1.38 \times 10^{-23}$ J/K o $8.617 \times 10^{-5}$ eV/K.
						\item $T$: Temperatura absoluta en Kelvin (K).
						\item $q$: Carga elemental del electrón. Es una constante física cuyo valor es $1,602\cdot10^{-19}\,\text{C}$. Se mide en Culombios ($C$).
						\item $N_A, N_D$: Concentración de aceptores y donadores.
						\item $n_i$: Concentración intrínseca.
					\end{itemize}

				\item \textbf{(5.16)} $\boxed{W = \left[ \frac{2\epsilon_0\epsilon_r V_{bi}}{q} \left( \frac{1}{N_A} + \frac{1}{N_D} \right) \right]^{1/2}}$ Ancho total de la zona de vaciamiento (o zona de deplexión) en equilibrio térmico.
					\begin{itemize}
						\item $W$: Ancho total de la región sin portadores libres. Se mide en metros (m) o centímetros (cm).
						\item $V_{bi}$: Potencial interno. Se mide en Voltios (V).
						\item $q$: Carga elemental del electrón. Es una constante física cuyo valor es $1,602\cdot10^{-19}\,\text{C}$. Se mide en Culombios ($C$).
						\item $\epsilon_{0}$: Permitividad eléctrica del vacío. Constante física cuyo valor es $8,854\cdot10^{-12}\,\text{F/m}$. Describe la resistencia que ofrece el vacío a la formación de un campo eléctrico.
						\item $N_A, N_D$: Concentración de aceptores y donadores.
					\end{itemize}

				\item \textbf{(5.17)} $\boxed{x_p = W \cdot \frac{N_D}{N_A + N_D}}$ Extensión de la zona de vaciamiento hacia el lado tipo p.
					\begin{itemize}
						\item $x_p$: Profundidad de penetración de la zona de vaciamiento en el lado p. Se mide en metros (m). Muestra que la zona penetra menos en el lado más dopado.
						\item $W$: Ancho total de la zona de vaciamiento.
						\item $N_D, N_A$: Concentración de donadores y aceptores.
					\end{itemize}

				\item \textbf{(5.18)} $\boxed{x_n = W \cdot \frac{N_A}{N_A + N_D}}$ Extensión de la zona de vaciamiento hacia el lado tipo n.
					\begin{itemize}
						\item $x_n$: Profundidad de penetración de la zona de vaciamiento en el lado n. Se mide en metros (m).
						\item $W$: Ancho total de la zona de vaciamiento.
						\item $N_D, N_A$: Concentración de donadores y aceptores.
					\end{itemize}

				\item \textbf{(5.19)} $\boxed{n(x = -x_p) = n_{p,0} \cdot e^{\frac{qV}{k_B T}}}$ Concentración de electrones (minoritarios) inyectados justo en el borde de la zona de vaciamiento del lado p al aplicar un voltaje externo $V$.
					\begin{itemize}
						\item $n(x = -x_p)$: Concentración de electrones en la posición $-x_p$ (borde de la región p). Se mide en $cm^{-3}$.
						\item $n_{p,0}$: Concentración de electrones minoritarios en equilibrio en la región p (lejos de la unión). Se mide en $cm^{-3}$.
						\item $q$: Carga elemental del electrón. Es una constante física cuyo valor es $1,602\cdot10^{-19}\,\text{C}$. Se mide en Culombios ($C$).
						\item $V$: Voltaje de polarización aplicado al diodo. Se mide en Voltios (V).
						\item $k_B$: Constante de Boltzmann. Su valor es aproximadamente $1,38 \cdot 10^{-23}\,\text{J/K}$ (o $8,617 \cdot 10^{-5}\,\text{eV/K}$).
						\item $T$: Temperatura absoluta del semiconductor. Se mide en Kelvin ($K$).
					\end{itemize}

				\item \textbf{(5.20)} $\boxed{p(x = x_n) = p_{n,0} \cdot e^{\frac{qV}{k_B T}}}$ Concentración de huecos (minoritarios) inyectados justo en el borde de la zona de vaciamiento del lado n.
					\begin{itemize}
						\item $p(x = x_n)$: Concentración de huecos en la posición $x_n$ (borde de la región n). Se mide en $cm^{-3}$.
						\item $p_{n,0}$: Concentración de huecos minoritarios en equilibrio en la región n. Se mide en $cm^{-3}$.
						\item $V$: Voltaje de polarización aplicado.
						\item $q$: Carga elemental del electrón. Es una constante física cuyo valor es $1,602\cdot10^{-19}\,\text{C}$. Se mide en Culombios ($C$).
						\item $k_B$: Constante de Boltzmann. Su valor es aproximadamente $1,38 \cdot 10^{-23}\,\text{J/K}$ (o $8,617 \cdot 10^{-5}\,\text{eV/K}$).
						\item $T$: Temperatura absoluta del semiconductor. Se mide en Kelvin ($K$).
					\end{itemize}

				\item \textbf{(5.21)} $\boxed{I(V) = I_s \left( e^{\frac{q \cdot V}{k_B T}} - 1 \right)}$ Ecuación del diodo ideal de Shockley. Relaciona la corriente que atraviesa el diodo con el voltaje aplicado.
					\begin{itemize}
						\item $I(V)$: Corriente a través del diodo en función del voltaje $V$. Se mide en Amperios (A).
						\item $I_s$: Corriente de saturación inversa. Es una constante muy pequeña que depende del material y la temperatura. Se mide en Amperios (A).
						\item $V$: Voltaje o diferencia de potencial aplicada al diodo. Positivo para polarización directa, negativo para inversa. Se mide en Voltios (V).
						\item $q$: Carga elemental del electrón. Es una constante física cuyo valor es $1,602\cdot10^{-19}\,\text{C}$. Se mide en Culombios ($C$).
						\item $k_B$: Constante de Boltzmann. Su valor es aproximadamente $1,38 \cdot 10^{-23}\,\text{J/K}$ (o $8,617 \cdot 10^{-5}\,\text{eV/K}$).
						\item $T$: Temperatura absoluta del semiconductor. Se mide en Kelvin ($K$).
					\end{itemize}
			\end{itemize}
\end{document}